\documentclass[12pt]{article}

\usepackage{sbc-template}

\usepackage[utf8]{inputenc}
\usepackage[T1]{fontenc}
\usepackage[brazil]{babel}

\usepackage{graphicx}
\usepackage{float}
\usepackage{booktabs}
\usepackage{caption}

\usepackage{amsmath,amssymb}

\usepackage{tikz}
\usetikzlibrary{positioning,arrows.meta,shapes,shadows}

\usepackage{xcolor}
\usepackage{listings}

\usepackage{url}
\usepackage{hyperref}
\hypersetup{hidelinks}

\lstset{
basicstyle=\ttfamily\small,
keywordstyle=\color{blue!70!black}\bfseries,
commentstyle=\color{gray},
stringstyle=\color{green!50!black},
showstringspaces=false,
breaklines=true,
numbers=left,
numberstyle=\tiny\color{gray},
frame=single,
language=Python,
tabsize=2,
breakatwhitespace=true
}

\sloppy

\title{Análise da Rede de Colaboração em Repositórios do GitHub Utilizando Teoria de Grafos}

\author{Davi Nunes Carvalho\inst{1}, Luiz Fernando Batista Moreira\inst{1},\\
  Josué Carlos Goulart dos Reis\inst{1}, João Victor Russo Marquito\inst{1}}

\address{Curso de Engenharia de Software --- Pontifícia Universidade Católica de Minas Gerais
  (PUC Minas)\\
  Belo Horizonte --- MG --- Brasil
  \email{\{dnunes,lfbmoreira,jcgereis,jmarquito\}@sga.pucminas.br}
}

\begin{document}

\maketitle

\begin{abstract}
  This work presents a tool for analyzing collaboration networks in GitHub
  repositories using Graph Theory. The system models interactions among
  collaborators (comments, reviews, merges, and issue closures) as directed
  weighted graphs. The implementation includes two graph representations
  (adjacency list and adjacency matrix), integration with the GitHub REST API,
  computation of 10+ network analysis metrics (PageRank, betweenness, closeness,
  eigenvector, clustering, assortativity, community detection, and bridging ties),
  and CSV export compatible with Gephi. The project was validated with 84 unit
  tests and executed over the \texttt{devlikeapro/waha} repository.
\end{abstract}

\begin{resumo}
  Este trabalho apresenta uma ferramenta para análise de redes de colaboração em
  repositórios do GitHub utilizando Teoria de Grafos. O sistema modela interações
  entre colaboradores (comentários, revisões, merges e fechamentos de issues)
  como grafos ponderados direcionados. A implementação inclui duas representações
  de grafos (lista de adjacência e matriz de adjacência), integração com a API
  REST do GitHub, cálculo de 10+ métricas de análise de rede (PageRank,
  betweenness, closeness, eigenvector, clustering, assortatividade, detecção de
  comunidades e bridging ties) e exportação CSV compatível com o Gephi. O projeto
  foi validado com 84 testes unitários e executado sobre o repositório
  \texttt{devlikeapro/waha}.
\end{resumo}

% ============================================================================
% 1. INTRODUÇÃO
% ============================================================================
\section{Introdução}\label{sec:intro}

Repositórios de software aberto no GitHub constituem redes sociotécnicas complexas onde colaboradores interagem através de múltiplos canais: comentários em \textit{issues}, discussões em \textit{pull requests}, revisões de código e decisões de integração. Essas interações formam uma estrutura subjacente que codifica padrões de influência, expertise e comunicação dentro do projeto.

A aplicação de Teoria de Grafos a essas redes permite quantificar e visualizar a dinâmica colaborativa. Representando usuários como vértices e interações como arestas ponderadas direcionadas, é possível identificar: (i) colaboradores centrais e influentes; (ii) estruturas de comunidades informais; (iii) usuários que atuam como conectores entre grupos isolados; (iv) padrões assortivos ou dissortativos de conexão. Essas informações são essenciais para compreender a governança e sustentabilidade de projetos open source.

\subsection{Objetivos}

\begin{enumerate}
  \item Implementar estruturas de dados de grafos em Python puro, sem bibliotecas especializadas, atendendo uma API completa com duas representações (\textit{adjacency list} e \textit{adjacency matrix})
  \item Modelar interações reais do GitHub como grafos ponderados direcionados, com acumulação automática de pesos em arestas paralelas
  \item Desenvolver um minerador robusto que coleta dados via API REST do GitHub, tratando autenticação, paginação e rate limits
  \item Implementar 10+ métricas de análise de redes complexas: centralidade de grau, betweenness, closeness, PageRank, eigenvector, clustering, assortatividade, densidade, detecção de comunidades e bridging ties
  \item Validar a ferramenta através de testes unitários e demonstrar exportação para Gephi
\end{enumerate}

% ============================================================================
% 2. MODELAGEM DO PROBLEMA
% ============================================================================
\section{Modelagem do Problema}\label{sec:modelagem}

\subsection{Justificativa da Escolha do Repositório}

O repositório \texttt{devlikeapro/waha} foi escolhido por possuir mais de 5.000 estrelas no GitHub, atendendo ao requisito de uma comunidade ativa e com grande volume de interações. O projeto WAHA (WhatsApp HTTP API) é uma solução open source que permite integrar aplicações ao WhatsApp por meio de APIs HTTP, sendo amplamente utilizada e mantida por diversos colaboradores.

Além da popularidade, o repositório apresenta uma quantidade significativa de issues, pull requests, revisões e comentários, fornecendo dados suficientes para modelar e analisar a rede de colaboração entre desenvolvedores. Essas características tornam o projeto adequado para a construção dos grafos de interação e para a aplicação das métricas de análise de redes propostas neste trabalho.

\subsection{Definição Formal}

O sistema constrói quatro grafos $G = (V, E)$ direcionados ponderados simples, onde:
\begin{align}
  V &= \{\text{logins únicos de usuários extraídos do repositório}\} \\
  E &\subseteq V \times V, \quad (u, v) \in E \Rightarrow u \neq v \text{ (sem self-loops)} \\
  w: E &\to \mathbb{R}^+ \quad \text{(peso acumulado de arestas)}
\end{align}

Os quatro grafos modelam diferentes dimensões de interação:
\begin{enumerate}
  \item \textbf{G1} --- comentários em issues e pull requests (peso 2 por interação)
  \item \textbf{G2} --- fechamentos de issues por outro usuário (peso não especificado no integrado)
  \item \textbf{G3} --- revisões/aprovações e merges de pull requests (pesos 4 e 5)
  \item \textbf{Grafo Integrado} --- fusão ponderada de G1, G2 e G3
\end{enumerate}

\subsection{Estratégia de Coleta de Dados}

O minerador \texttt{GithubMiner} segue: (i) autenticação via token em \texttt{.env}; (ii) paginação de issues (\texttt{per\_page=100}, 5 páginas); (iii) para cada item: coleta comentários, revisões, detalhes de merge; (iv) tratamento robusto de erros (ignora status $\neq$ 200).

\subsection{Transformação de Interações em Arestas}

Cada interação é convertida em aresta direcionada com acumulação de pesos:
\begin{itemize}
  \item \textbf{Comentário:} usuário $c$ comenta issue/PR aberta por $a \Rightarrow$ aresta $c \to a$ recebe $+2$ (G1 e integrado)
  \item \textbf{Abertura comentada:} se issue (não PR) recebe comentário de $c \neq a \Rightarrow$ aresta $a \to c$ recebe $+3$ (integrado; sinaliza envolvimento da comunidade)
  \item \textbf{Revisão/aprovação:} usuário $r$ revisa PR de $a \Rightarrow$ aresta $r \to a$ recebe $+4$ (G3 e integrado)
  \item \textbf{Merge:} usuário $m$ mergua PR de $a \Rightarrow$ aresta $m \to a$ recebe $+5$ (G3 e integrado)
  \item \textbf{Acumulação:} múltiplas interações do mesmo par $(u, v)$ somam seus pesos (e.g., 3 comentários = peso 6)
\end{itemize}

\subsection{Esquema de Pesos}

\begin{table}[H]
  \centering
  \caption{Pesos das interações no grafo integrado. Hierarquia reflete relevância para colaboração técnica.}
  \label{tab:pesos}
  \begin{tabular}{l|c|p{5.5cm}}
    \toprule
    \textbf{Interação} & \textbf{Peso} & \textbf{Semântica} \\
    \midrule
    Comentário em issue/PR & 2 & Participação leve (opinião, sugestão) \\
    Abertura comentada & 3 & Comunidade engajada na proposta \\
    Revisão/aprovação de PR & 4 & Análise técnica profunda \\
    Merge de PR & 5 & Decisão executiva de integração \\
    \bottomrule
  \end{tabular}
\end{table}

A hierarquia de pesos reflete a tese de que revisões e merges requerem expertise maior que comentários: um revisor avalia correção técnica; um merguer assume responsabilidade de integração. Comentários são participação essencial mas menos comprometida.

% ============================================================================
% 3. MODELAGEM DO SISTEMA
% ============================================================================
\section{Modelagem do Sistema}\label{sec:modelagem_sistema}

Esta seção apresenta os artefatos UML que descrevem a arquitetura estática, comportamental e os fluxos de execução da ferramenta.

\subsection{Diagrama de Classes}

O sistema organiza sua arquitetura em três principais responsabilidades. A classe abstrata \texttt{AbstractGraph} define uma interface comum para operações sobre grafos, sendo implementada pelas classes concretas \texttt{AdjacencyListGraph} e \texttt{AdjacencyMatrixGraph}. A primeira utiliza uma estrutura baseada em lista de adjacência \texttt{List[Dict[int, float]]}, sendo adequada para grafos esparsos, enquanto a segunda utiliza uma matriz de adjacência $N \times N$ com valores \texttt{float|None}, permitindo acesso direto às arestas. A classe \texttt{GithubMiner} é responsável pela coleta dos dados da API e pela construção das estruturas de grafo. Já a classe \texttt{GraphAnalytics} concentra os algoritmos de análise, incluindo métricas de centralidade, agrupamento e relacionamento entre os colaboradores.

\textbf{Estrutura UML:} Classe abstrata \texttt{AbstractGraph} com 20+ métodos define interface comum. Implementações concretas \texttt{AdjacencyListGraph} e \texttt{AdjacencyMatrixGraph} fornecem representações complementares. \texttt{GithubMiner} coleta via API, \texttt{GraphAnalytics} executa algoritmos.

\subsection{Diagrama de Componentes}

Arquitetura modular em 5 componentes: (i) \textit{Main Application} --- orquestra fluxo; (ii) \textit{GitHub Miner} --- coleta API; (iii) \textit{Graph Storage} --- instancia grafos; (iv) \textit{Analytics Engine} --- métricas; (v) \textit{Export Module} --- saída CSV/relatórios.

\subsection{Fluxo de Mineração}

(1) Aplicação invoca \texttt{fetch\_data(pages=5)}; (2) minerador requisita \texttt{/repos/owner/repo/issues}; (3) para cada item coleta comentários, revisões, merges; (4) acumula pesos em dicionários por tipo; (5) \texttt{build\_graphs()} mapeia usuários a IDs, instancia 4 grafos \texttt{AdjacencyListGraph}; (6) exporta para CSV (Source, Target, Weight).

% ============================================================================
% 4. ARQUITETURA DA FERRAMENTA
% ============================================================================
\section{Arquitetura da Ferramenta}\label{sec:arquitetura}

A ferramenta está organizada em 3 camadas:

\begin{figure}[H]
  \centering
  \begin{tikzpicture}[
    node distance=0.9cm,
    every node/.style={draw, rounded corners, minimum width=6cm, align=center, font=\small},
    app/.style={fill=yellow!20},
    analytics/.style={fill=orange!20},
    miner/.style={fill=blue!20},
    core/.style={fill=green!20},
  ]
    \node[app] (app) {\textbf{Aplicação}\\main.py, src/app.py};
    \node[analytics, below=of app] (analytics) {\textbf{Análise}\\GraphAnalytics};
    \node[miner, below=of analytics] (miner) {\textbf{Mineração}\\GithubMiner};
    \node[core, below=of miner] (core) {\textbf{Núcleo de Grafos}\\AbstractGraph, AdjacencyListGraph, AdjacencyMatrixGraph};

    \draw[-{Stealth[length=2mm]}, thick] (app) -- (analytics);
    \draw[-{Stealth[length=2mm]}, thick] (app) -- (miner);
    \draw[-{Stealth[length=2mm]}, thick] (miner) -- (core);
    \draw[-{Stealth[length=2mm]}, thick] (analytics) -- (core);
  \end{tikzpicture}
  \caption{Arquitetura em camadas da ferramenta.}
  \label{fig:arquitetura}
\end{figure}

\subsection{Camada 1: Núcleo de Grafos}

\texttt{AbstractGraph} define 20+ métodos: gerenciamento (getVertexCount, getEdgeCount), arestas (addEdge, removeEdge, hasEdge, pesos), graus (in/out-degree), relações (sucessor, predecessor, convergência), propriedades (conectividade, completude), exportação (exportToGEPHI em CSV).

\texttt{AdjacencyListGraph} usa listas de dicionários ($O(E)$ espaço, eficiente para redes esparsas). \texttt{AdjacencyMatrixGraph} usa matriz $N \times N$ (acesso $O(1)$ a arestas, cache de in-degrees).

\subsection{Camada 2: Mineração --- \texttt{GithubMiner}}

Responsável por toda interação com a API REST do GitHub. Recebe token e repositório via \texttt{.env}.

\begin{itemize}
  \item \texttt{fetch\_data(pages)}: percorre issues e PRs, coleta comentários e revisões, acumula pesos em dicionários por tipo de grafo (\texttt{edges\_g1}, \texttt{edges\_g2}, \texttt{edges\_g3}, \texttt{edges\_integrated})
  \item \texttt{build\_graphs()}: mapeia logins para IDs sequenciais e instancia os 4 grafos como \texttt{AdjacencyListGraph}
\end{itemize}

\subsection{Camada 3: Análise --- \texttt{GraphAnalytics}}

Recebe um \texttt{AbstractGraph} e implementa todos os algoritmos em Python puro, sem bibliotecas externas de grafos.

\subsection{Restrições Atendidas}

Sem bibliotecas externas (NetworkX); sem self-loops (verificação $u \neq v$); sem múltiplas arestas (acumulação de pesos); idempotência garantida; validação robusta de índices.

\subsection{Testes Unitários}

84 testes cobrindo: (i) API de grafos (ambas implementações): arestas, graus, relações, propriedades; (ii) mineração: mock da API, parsing, acumulação de pesos; (iii) casos extremos (grafos vazios, índices inválidos). Todos passam.

\subsection{Aplicação de Demonstração}

\texttt{src/app.py} demonstra a API: criação de grafos, adição de arestas, consulta de graus, verificação de conectividade, exportação para Gephi.

% ============================================================================
% 5. MÉTRICAS E ANÁLISE
% ============================================================================
\section{Métricas e Análise}\label{sec:metricas}

Todas as métricas são implementadas na classe \texttt{GraphAnalytics} em Python puro.

\subsection{Centralidade de Grau}

\[
C_D(v) = \frac{\text{in-degree}(v) + \text{out-degree}(v)}{2(n-1)}
\]

Mede quantas conexões diretas cada colaborador tem. In-degree = usuários que interagiram \textit{com} $v$; out-degree = usuários que $v$ interagiu \textit{com}.

\subsection{Betweenness Centrality}

\[
C_B(v) = \sum_{s \neq v \neq t} \frac{\sigma_{st}(v)}{\sigma_{st}}
\]

Mede quantos caminhos mais curtos passam por $v$. Valor alto = "conector" ou "ponte" entre grupos. Implementado via algoritmo de Brandes $O(VE)$.

\subsection{Closeness Centrality}

\[
C_C(v) = \frac{n-1}{\sum_{t \neq v} d(v,t)}
\]

Mede proximidade a todos os outros. Valor alto = colaborador próximo a todos (poucos hops). Identifica quem tem acesso rápido à informação.

\subsection{PageRank}

\[
\text{PR}(v) = \frac{1-d}{n} + d \sum_{u \to v} \frac{\text{PR}(u)}{\text{out-degree}(u)}
\]

Damping factor $d = 0{,}85$. Implementado via power iteration. Mede influência ponderada pela importância dos conectados.

\subsection{Eigenvector Centrality}

\[
x_v = \lambda^{-1} \sum_{u \to v} x_u
\]

Mede influência via importância recursiva dos vizinhos. Resolvido via power iteration.

\subsection{Clustering, Densidade, Assortatividade e Comunidades}

\textbf{Clustering:} $C(v) = \frac{\text{arestas entre vizinhos}}{\binom{k_v}{2}}$ mede triângulos locais. \textbf{Densidade:} $\rho = \frac{|E|}{|V|(|V|-1)}$ = proporção de conexões. \textbf{Assortatividade:} correlação de Pearson sobre graus; positiva = preferência por pares similares. \textbf{Comunidades:} Louvain detecta subgrupos. \textbf{Bridging ties:} usuários conectando comunidades distintas.

% ============================================================================
% 6. RESULTADOS E DISCUSSÃO
% ============================================================================
\section{Resultados e Discussão}\label{sec:resultados}

\subsection{Estatísticas Gerais}

A ferramenta foi executada sobre o repositório \texttt{devlikeapro/waha} (repositório de automação WhatsApp), coletando 5 páginas de issues e pull requests (máximo 500 itens). O dataset compreende o histórico completo de interações do repositório até a data de execução.

\begin{table}[H]
  \centering
  \caption{Estatísticas da rede analisada (repositório devlikeapro/waha).}
  \label{tab:stats}
  \begin{tabular}{l|r}
    \toprule
    \textbf{Métrica} & \textbf{Valor} \\
    \midrule
    Número de vértices (colaboradores únicos) & 481 \\
    Número de arestas --- Grafo Integrado & 1409 \\
    Densidade da rede & 0,006103 \\
    Grafo fracamente conectado & Não \\
    Clustering médio & 0,148564 \\
    Assortatividade & $-0{,}249$ (dissortativo) \\
    Comunidades detectadas (Louvain) & 135 \\
    \bottomrule
  \end{tabular}
\end{table}

\subsection{Análise de Estrutura}

\begin{table}[H]
  \centering
  \caption{Métricas estruturais da rede (repositório devlikeapro/waha).}
  \label{tab:estrutura}
  \begin{tabular}{l|r|p{6cm}}
    \toprule
    \textbf{Métrica} & \textbf{Valor} & \textbf{Interpretação} \\
    \midrule
    Densidade & 0,006044 & Rede muito esparsa; típica de open source com periferia dispersa \\
    Clustering médio & 0,1445 & Triângulos raros; baixa tendência de formação de cliques \\
    Assortatividade & $-0,250$ & Padrão dissortativo: hubs conectam-se com periféricos \\
    Comunidades (Louvain) & 137 & Muitos isolados; estrutura fragmentada \\
    \bottomrule
  \end{tabular}
\end{table}

\subsection{Centralidade de Grau --- Top Colaboradores}

\begin{table}[H]
  \centering
  \caption{Graus (in/out) dos 5 principais colaboradores.}
  \label{tab:graus}
  \begin{tabular}{l|r|r|r}
    \toprule
    \textbf{Usuário} & \textbf{In-Degree} & \textbf{Out-Degree} & \textbf{Total} \\
    \midrule
    devlikepro           & 185 & 192 & 377 \\
    bergpinheiro         & 86  & 94  & 180 \\
    vicentemarciot27     & 22  & 18  & 40  \\
    github-actions[bot]  & 25  & 28  & 53  \\
    OrionZap             & 26  & 15  & 41  \\
    \bottomrule
  \end{tabular}
\end{table}

\subsection{Betweenness Centrality --- Conectores da Rede}

\begin{table}[H]
  \centering
  \caption{Top 5 colaboradores por betweenness centrality (pontes da rede).}
  \label{tab:betweenness}
  \begin{tabular}{l|r}
    \toprule
    \textbf{Usuário} & \textbf{Betweenness} \\
    \midrule
    devlikepro         & 0,6078 \\
    bergpinheiro       & 0,1641 \\
    vicentemarciot27   & 0,0570 \\
    github-actions[bot] & 0,0457 \\
    ehussain           & 0,0414 \\
    \bottomrule
  \end{tabular}
\end{table}

\textbf{Devlikepro} (0,608) é a ponte dominante: 60,8\% de todos caminhos mais curtos passam por ela. \textbf{Bergpinheiro} (0,164) é a segunda; concentra 16,4\%.

\subsection{Closeness Centrality --- Proximidade}

\begin{table}[H]
  \centering
  \caption{Top 5 usuários isolados e devlikepro por closeness centrality.}
  \label{tab:closeness}
  \begin{tabular}{l|r}
    \toprule
    \textbf{Usuário} & \textbf{Closeness} \\
    \midrule
    KempesBarbosa      & 1,0000 \\
    NatanaelRibF       & 1,0000 \\
    VikthorPenaloz     & 1,0000 \\
    nicolasvahidzein   & 1,0000 \\
    reeshan            & 1,0000 \\
    devlikepro         & 0,6017 \\
    \bottomrule
  \end{tabular}
\end{table}

Closeness 1,0 = usuários isolados (distância infinita aos outros = inf normalizado a 1). Devlikepro (0,6017) = mais próxima a todos entre os ativos.

\subsection{PageRank e Eigenvector Centrality --- Influência}

\begin{table}[H]
  \centering
  \caption{Top 5 colaboradores: PageRank vs. Eigenvector Centrality.}
  \label{tab:centrality}
  \begin{tabular}{l|r|r}
    \toprule
    \textbf{Usuário} & \textbf{PageRank} & \textbf{Eigenvector} \\
    \midrule
    devlikepro         & 0,1171 & 0,5895 \\
    bergpinheiro       & 0,0508 & 0,2947 \\
    vicentemarciot27   & 0,0175 & --- \\
    OrionZap           & 0,0139 & 0,1441 \\
    github-actions[bot] & 0,0134 & 0,1581 \\
    \bottomrule
  \end{tabular}
\end{table}

PageRank e Eigenvector correlacionam: devlikepro 11,7\% (PageRank) e 58,9\% (Eigenvector) = principal influenciador.

\subsection{Clustering Coefficient}

Coeficiente médio: 0,1445. Top 10 com máximo local (1,0): usuários formando triângulos fechados. Exemplo: \textit{Biasolis}, \textit{DiegoAugustusCunha}, etc. formam subgrupos coesos.

\subsection{Detecção de Comunidades}

\begin{table}[H]
  \centering
  \caption{Distribuição de comunidades (top 5 maiores).}
  \label{tab:comunidades}
  \begin{tabular}{l|r}
    \toprule
    \textbf{Comunidade} & \textbf{Tamanho} \\
    \midrule
    Comunidade 15 & 67 membros \\
    Comunidade 19 & 33 membros \\
    Comunidade 73 & 22 membros \\
    Comunidade 9  & 14 membros \\
    Comunidade 34 & 12 membros \\
    \bottomrule
  \end{tabular}
\end{table}

137 comunidades detectadas. Maior (67 membros) = ~14\% da rede. Muitas comunidades pequenas ($\leq 3$ membros).

\subsection{Bridging Ties --- Conectores entre Comunidades}

\begin{table}[H]
  \centering
  \caption{Top 5 usuários que conectam múltiplas comunidades.}
  \label{tab:bridging}
  \begin{tabular}{l|r}
    \toprule
    \textbf{Usuário} & \textbf{Comunidades conectadas} \\
    \midrule
    devlikepro       & 79 \\
    bergpinheiro     & 37 \\
    github-actions[bot] & 20 \\
    OrionZap         & 17 \\
    matteus-siqueira & 13 \\
    \bottomrule
  \end{tabular}
\end{table}

Devlikepro conecta 79 de 137 comunidades (57,7\%); essencial para coesão global. Bergpinheiro conecta 37 (27\%).

\subsection{Exportação para Gephi}

CSV com 1409 arestas (Source, Target, Weight). Gephi visualiza: núcleo central (top-5 com maior in-degree), periferia esparsa, comunidades por cor modularity, tamanho de nó por degree.

% ============================================================================
% 7. RESPONSABILIDADES DOS INTEGRANTES
% ============================================================================
\section{Responsabilidades dos Integrantes}\label{sec:responsabilidades}

A seguir, resumo das contribuições de cada integrante do grupo:

\begin{itemize}
  \item \textbf{Davi Nunes Carvalho --- Mineração e Núcleo de Grafos:} responsável pelo minerador \texttt{GithubMiner} (coleta de issues, PRs, comentários, revisões e merges via API REST do GitHub) e pelas estruturas de dados de grafos (\texttt{AbstractGraph} e \texttt{AdjacencyListGraph}). Implementou também a exportação para o Gephi e a integração com a API.

  \item \textbf{Luiz Fernando Batista Moreira --- Análise de Redes:} responsável pela classe \texttt{GraphAnalytics} e pela implementação, em Python puro, dos algoritmos de análise: PageRank, betweenness, closeness, eigenvector centrality e detecção de comunidades (Louvain).

  \item \textbf{Josué Carlos Goulart dos Reis --- Testes e Validação:} responsável pelos 84 testes unitários, pelo mock da API para testes, pela validação da mineração e pela documentação técnica do projeto.

  \item \textbf{João Victor Russo Marquito --- Modelagem e Documentação:} responsável pelos diagramas UML, pela redação do relatório em \LaTeX{}, pela definição formal da modelagem do grafo e pela interpretação dos resultados obtidos.
\end{itemize}

% ============================================================================
% 8. CONCLUSÃO
% ============================================================================
\section{Conclusão}\label{sec:conclusao}

Este trabalho implementou com sucesso uma ferramenta de análise de redes de colaboração em GitHub atendendo todos os requisitos das 3 etapas:

\textbf{Etapa 1:} Modelagem com 4 grafos (G1, G2, G3, integrado), estratégia de coleta de 5 tipos de interação, esquema de pesos refletindo importância técnica.

\textbf{Etapa 2:} Implementação de estruturas de dados em Python puro (\texttt{AbstractGraph}, \texttt{AdjacencyListGraph}, \texttt{AdjacencyMatrixGraph}), API completa (20+ métodos), 84 testes unitários, exportação Gephi, minerador robusto com tratamento de erros.

\textbf{Etapa 3:} Análise via 10+ métricas: grau, betweenness, closeness, PageRank, eigenvector, clustering, densidade, assortatividade, detecção de comunidades (Louvain), bridging ties. Resultados validados e interpretados sobre repositório com 481 usuários e 1409 arestas.

\textbf{Principais achados:} Devlikepro é figura central (PR 0,117, betweenness 0,608, conecta 79 comunidades). Rede dissortativa (hub-and-spoke) com estrutura fragmentada (137 comunidades). Densidade 0,006 = típica de open source com periferia dispersa.

\begin{thebibliography}{9}

\bibitem{brandes2001}
  Ulrik Brandes.
  \textit{A Faster Algorithm for Betweenness Centrality}.
  Journal of Mathematical Sociology, 25(2):163--177, 2001.

\bibitem{page1999}
  Lawrence Page, Sergey Brin, Rajeev Motwani, Terry Winograd.
  \textit{The PageRank Citation Ranking: Bringing Order to the Web}.
  Stanford InfoLab, 1999.

\bibitem{newman2010}
  M.\ E.\ J.\ Newman.
  \textit{Networks: An Introduction}.
  Oxford University Press, 2010.

\bibitem{gephi}
  Mathieu Bastian, Sebastien Heymann, Mathieu Jacomy.
  \textit{Gephi: an open source software for exploring and manipulating networks}.
  ICWSM, 2009.

\bibitem{github_api}
  GitHub Documentation.
  \textit{GitHub REST API}.
  \url{https://docs.github.com/en/rest}, 2024.

\end{thebibliography}

\end{document}
